\begin{abstract}
While large foundation models are now widespread, the computational demands of training them from the ground up are immense. Therefore, methods for efficiently adapting these sophisticated models to new tasks have become increasingly crucial. This work focuses on a robust finetuning approach termed Orthogonal Finetuning (OFT), which adapts pretrained models for specific tasks. Although OFT shows notable transferability, its reliance on high-dimensional orthogonal matrices entails a substantial number of trainable parameters. To remedy this, we analyze OFT through the lens of information transfer, deriving several core criteria necessary for parameter efficiency. Motivated by the efficient communication structure in the Cooley-Tukey fast Fourier transform, we devise an orthogonal parameterization scheme based on butterfly structures. Incorporating this into OFT, we introduce a novel approach dubbed Orthogonal Butterfly (BOFT), which allows for greater parameter-efficiency and generalizes OFT as a special instance. We conduct comprehensive experiments on a range of large-scale models—including vision transformers, language models, and text-to-image diffusion architectures—demonstrating the efficacy of our proposed method for various vision and language downstream tasks.
\end{abstract}

\section{Introduction}

Advances exemplified by models such as ChatGPT~\cite{brown2020language,chen2021evaluating} and Stable Diffusion~\cite{rombach2022high} reveal the extensive generalization capacities of large foundation models. Accompanying these performance leaps is a steep rise in model parameter counts 

(\emph{e.g.}, GPT-3 possesses approximately 175B parameters). Consequently, the practicalities of training foundation models from the beginning have become a significant obstacle for most researchers. For continued progress in the area, it is essential to develop techniques that can tailor pretrained models to new tasks without resorting to full retraining. This necessity has led to three main adaptation paradigms: \emph{model finetuning}~\cite{hulora2022,zaken2022bitfit,guo2020parameter,raffel2020exploring,qiu2023controlling,zhang2022adaptive,chavan2023one}, which selectively updates a subset of the original model's parameters without altering the architecture; \emph{adapter tuning}~\cite{rebuffi2017learning,houlsby2019parameter,pfeiffer2020adapterfusion,lin2020exploring,he2021towards}, involving the integration of additional learnable modules; and \emph{prompt tuning}~\cite{li2021prefix,lester2021power}, where trainable prompt tokens are prepended to the input sequence. Among these strategies, model finetuning is particularly attractive due to its conceptual simplicity, effectiveness, and the advantage of not increasing inference time.

The core idea underlying model finetuning is to maintain similarity between the pretrained and finetuned models according to specific metrics, thereby retaining the knowledge acquired during pretraining. In practice, existing finetuning approaches frequently employ a reduced learning rate, an informal technique aimed at limiting the difference between the pretrained and updated models. Considering the intrinsic structure of weight matrices, a more systematic strategy seeks to conserve the relational properties within these matrices, specifically the angles between neuron pairs. Building on this perspective, Orthogonal Finetuning~(OFT)~\cite{qiu2023controlling} was introduced, where all neurons in a layer are jointly transformed using a common orthogonal matrix, thus guaranteeing that neuron pairwise angles remain intact over the course of finetuning. While OFT has yielded notable improvements in both generalizability and optimization for text-to-image diffusion model finetuning~\cite{qiu2023controlling}, it is challenged by the substantial parameter count that arises from the large size of orthogonal matrices. To mitigate this, OFT adopts a block-diagonal formulation, cutting down on the number of learnable parameters. Nevertheless, this approach imposes certain limitations: the orthogonal matrix’s fixed sparse structure and the application of independent transformations within separate blocks can hinder expressivity and introduce potentially problematic inductive biases (\emph{e.g.}, such block-diagonal orthogonal matrices are less expressive and unable to represent standard linear transformations). Additionally, identifying an optimal block-wise sparsity pattern remains an unresolved challenge.

A central challenge is to construct a dense orthogonal matrix in a manner that is highly parameter-efficient. Conventional wisdom suggests that representing a dense orthogonal matrix in $d$ dimensions demands $\mathcal{O}(d^2)$ parameters, rendering dense orthogonalization seemingly prohibitive. We address this by innovatively assembling a dense orthogonal matrix through the multiplication of several factorized sparse matrices. This strategy stems from the observation that a reduction in the number of trainable parameters is achievable by accepting increased computational complexity in exchange for decreased memory footprint. In our formulation, the orthogonal matrix is implicitly realized by the sequential multiplication of sparse matrices, necessitating repeated multiplication steps during each iteration of training. To conceptualize this, we model the synthesis of a dense orthogonal matrix as an instance of information propagation. In this lens, constructing a dense orthogonal matrix via successive sparse matrices mirrors the process of transmitting information across a grid-structured network. This interpretation empowers us to devise a broad class of sparse matrix factorization schemes that restrict trainable parameter count while retaining sufficient expressive power to generate fully dense matrices.

For enhanced parameter efficiency within this information transmission paradigm, we take cues from butterfly structures characteristic of the Cooley-Tukey fast Fourier transform algorithm~\cite{cooley1965algorithm}, where efficient information exchange is facilitated among nodes~\cite{klasing1994broadcasting}. The butterfly network inherent in the Cooley-Tukey algorithm inspires a corresponding sparse matrix factorization technique, known as butterfly factorization. With butterfly factorization, a $d$-dimensional dense matrix can be represented as a product of $\mathcal{O}(\log d)$ sparse matrices, with each matrix containing only $\mathcal{O}(d)$ nonzero elements. By ensuring the orthogonality of each constituent sparse matrix, this factorization yields a parameterization that requires only $\mathcal{O}(d\log d)$ parameters, a substantial reduction compared to standard approaches. Capitalizing on this efficient orthogonal representation, we introduce Orthogonal Butterfly~(BOFT), a new finetuning approach that is both parameter-efficient and flexible. BOFT generalizes OFT by subsuming it as a particular instance, establishing a comprehensive framework for orthogonal finetuning. Both the block-diagonal and butterfly patterns share the key attribute of sparsity, which serves to decrease the number of learnable parameters in orthogonal matrices. Notably, there is a compelling parallel between our method and the low-rank designs of LoRA~\cite{hulora2022}; low-rank and sparse matrices represent two principal forms of structured matrices~\cite{candes2011robust} commonly utilized to promote parameter efficiency.

In contrast to the block-diagonal parameterization adopted in OFT, which balances model expressiveness with structural regularity, BOFT leverages the butterfly structure to create a more continuous transition between full orthogonal matrices (offering expressivity) and identity matrices (imposing regularity). This strategy allows us to identify a more compact hypothesis class within the space of orthogonal matrices for downstream applications. The butterfly pattern, being a foundational building block in widely-used fast linear operators—such as the discrete Fourier transform and the discrete cosine transform—suggests that our structured parameterization method inherently introduces an inductive bias, which we anticipate will improve generalization and guard against overfitting. This perspective is informed by the fact that the butterfly construction can reproduce numerous well-known linear operators, a feat unattainable by OFT’s block-diagonal design. Our key contributions can be summarized as follows:

\begin{itemize}[leftmargin=*,nosep]
    \item We tackle the challenge of parameter efficiency in orthogonal finetuning by introducing a novel framework grounded in information transmission theory. Specifically, we recast the challenge of constructing a parameter-efficient dense orthogonal matrix as an information propagation problem on a grid-like graph structure.
    \item Drawing inspiration from the butterfly schemes integral to the Cooley-Tukey algorithm, we put forward Orthogonal Butterfly—a new approach for parameter-efficient orthogonal finetuning—formulated within the context of our information transmission framework.
    \item We provide several theoretical explanations elucidating how BOFT achieves drastic reductions in the number of tunable parameters while preserving both expressive power and training stability. Due to its matrix decomposition structure, BOFT also facilitates a rich space of weight interpolation mechanisms (see Figure~\ref{fig:boft_interpolation}).
    \item For the first time, we expand orthogonal finetuning beyond controllable text-to-image generation~\cite{qiu2023controlling}, evaluating its utility in a broad array of tasks. We apply BOFT to various adaptation problems in computer vision and natural language processing. Our experiments demonstrate that BOFT surpasses the latest parameter-efficient finetuning methods by a significant margin, affirming its advantages in both parameter usage and generalization performance.
\end{itemize}

\textbf{Parameter-efficient finetuning (PEFT)}. As the scale and capability of foundation models (\emph{e.g.}, \cite{brown2020language,radford2021learning,rombach2022high,kirillov2023segment}) expand, there has been substantial attention on methods that enable finetuning for downstream applications with minimal additional parameters~\cite{treviso2023efficient,ding2023parameter,lialin2023scaling}. Numerous PEFT strategies have been proposed~\cite{houlsby2019parameter,aghajanyan2020intrinsic,hulora2022,edalati2022krona,wang2022adamix,gheini2021cross,zaken2022bitfit,guo2020parameter,sung2021training,ansell2022composable,lester2021power,li2021prefix,vu2022spot,he2021towards,mao2021unipelt,karimi2021compacter,liu2022few,sung2022lst,chen2022parameter,jia2022visual,chen2022adaptformer,zhang2022neural,jie2023fact,lian2022scaling,luo2023towards}, among which approaches based on reparameterization~\cite{aghajanyan2020intrinsic,hulora2022,edalati2022krona} are particularly aligned with this work. In particular, LoRA~\cite{hulora2022} adapts pretrained weights by incorporating the product of two low-rank matrices, which yields strong outcomes on a range of natural language tasks. Since LoRA employs a uniform rank for these additional matrices, subsequent techniques~\cite{zhang2022adaptive,valipour2022dylora,zhang2023increlora} have introduced dynamic rank adjustment at different layers, ensuring more optimal parameter allocation. Due to the practical advantages, such low-rank modifications to weight matrices have been widely adopted~\cite{dettmers2023qlora,zi2023delta,chavan2023one}. Drawing inspiration from the use of hyperspherical energy for understanding generalization~\cite{liu2018learning,liu2021orthogonal}, \cite{qiu2023controlling} developed orthogonal finetuning (OFT), a distinct approach that modifies neurons in the same layer via orthogonal transformations. OFT demonstrates both better generalization and more stable training in comparison to LoRA, especially in finetuning text-to-image diffusion models. However, a notable limitation is that OFT typically introduces a greater number of trainable parameters relative to LoRA. Therefore, enhancing the parameter efficiency of OFT remains an important objective. Furthermore, the adaptability of OFT to a broader class of adaptation problems—beyond just controlling text-to-image diffusion models~\cite{qiu2023controlling}—is still unclear. BOFT addresses this by leveraging butterfly factorization to reduce the parameter count required for OFT, thereby facilitating the application of orthogonal finetuning to general adaptation tasks.

\textbf{Butterfly structures}. The radix-2 Cooley-Tukey algorithm decomposes an $N$-point discrete Fourier transform into two smaller $\frac{N}{2}$-point transforms at each stage, which gives rise to the so-called butterfly structure—representable as a product of several sparse matrices (collectively termed a butterfly matrix). Butterfly matrices have found utility in various contexts: parameterizing orthogonal matrices to circumvent the need for pivoting in Gaussian elimination and boost computational efficiency~\cite{parker1995random}; stabilizing the optimization of recurrent neural networks~\cite{jing2017tunable}; and facilitating kernel approximation~\cite{munkhoeva2018quadrature}. Previous work~\cite{dao2019learning,dao2019kaleidoscope} has explored learning rapid linear mappings via butterfly parameterizations, while butterfly matrices have also enabled scalable neural network training~\cite{chen2022pixelated,dao2022monarch}. Moreover, these structures contribute to fast matrix-vector multiplication~\cite{michielssen1996multilevel,o2010algorithm}, efficient matrix approximations~\cite{li2015butterfly}, and enhanced network communications~\cite{klasing1994broadcasting,li2003linear,rajasingh2016transmission}. Our approach diverges from these existing applications by employing butterfly structures specifically to augment the parameter efficiency of OFT.

\section{An Information Transmission View on Orthogonal Finetuning}

We begin by reviewing foundational aspects of OFT. In this approach, the finetuning process involves reparameterizing the new weight matrix as a product between a learnable orthogonal matrix and the fixed pretrained weight matrix. This contrasts with LoRA, which updates the weights by adding a low-rank matrix; in contrast, OFT employs a multiplicative, orthogonal matrix to modulate the pretrained weights. LoRA achieves parameter-efficiency by leveraging low-rank matrices, whereas classical OFT relies on a block-diagonal orthogonal structure~\cite{qiu2023controlling}—the smaller the diagonal blocks, the greater the efficiency of OFT. Figure~\ref{fig:comp_lora} offers a visual summary of these differences. The central rationale behind utilizing orthogonal transformations in finetuning is the retention of the angular relationships between neuron representations~\cite{liu2018learning,liu2021orthogonal,qiu2023controlling}, which helps preserve the semantic knowledge acquired during pretraining.

More specifically, OFT introduces an orthogonal matrix $\bm{R}\in\mathbb{R}^{d\times d}$ for a given pretrained linear layer $\bm{W}^0\in\mathbb{R}^{d\times n}$, modifying the standard forward computation from $\bm{z} = (\bm{W}^0)^\top\bm{x}$ to $\bm{z} = (\bm{R}\bm{W}^0)^\top\bm{x}$, with $\bm{x} \in \mathbb{R}^{d}$ as the input and $\bm{z} \in \mathbb{R}^{n}$ as the output. For initialization, $\bm{R}$ is set as the identity matrix to ensure zero-initialization. During finetuning, to guarantee that $\bm{R}$ remains orthogonal, Cayley parameterization is adopted, as in \cite{qiu2023controlling,liu2021orthogonal}: $\bm{R} = (\bm{I}+\bm{Q})(\bm{I}-\bm{Q})^{-1}$, with $\bm{Q}$ a skew-symmetric matrix satisfying $\bm{Q} = -\bm{Q}^\top$.

To improve parameter-efficiency, OFT employs a block-diagonal construction, expressing $\bm{R}$ as $\text{diag}(\bm{R}_1, \bm{R}_2, \cdots, \bm{R}_r)$, where each $\bm{R}_i\in\mathbb{R}^{b\times b}$ is a smaller orthogonal matrix and $br=d$. This block-diagonal structure achieves efficiency at the cost of imposing a grouping scheme: the dimensions of each neuron (i.e., columns in $\bm{W}^0$) are split into $r$ disjoint groups, with each group being transformed by a separate orthogonal matrix. Despite the empirical success of this design, there is no principled justification for partitioning a neuron's dimensions in this way, leading to an argument in favor of dense orthogonal matrices. This prompts the following question: \emph{Is it possible to design a dense orthogonal matrix that retains parameter-efficiency?}

In response to this problem, we propose constructing a dense orthogonal matrix as a product of several sparse orthogonal matrices. To unify and clarify the analysis of sparsity patterns in orthogonal matrix decompositions, we reinterpret the task of forming a dense orthogonal matrix in OFT as an instance of information transmission. In particular, expressing a dense matrix $\bm{R}\in\mathbb{R}^{d\times d}$ as a multiplication of $m$ square matrices, $\thickmuskip=2mu \medmuskip=2mu \bm{R}=\bm{B}_m\bm{B}_{m-1}\cdots\bm{B}_1$, is analogous to the passage of information across a $d\times (m+1)$ grid of nodes, as depicted in Figure~\ref{fig:info_trans}. This viewpoint is motivated by recognizing that a $d$-dimensional dense matrix describes all-to-all connectivity from a group of $d$ input nodes to $d$ output nodes. Each entry $R_{ij}$ of $\bm{R}$ being nonzero reflects the ability to transfer information from the $j$-th source node to the $i$-th destination node, whereas a zero entry denotes a blocked path. Hence, decomposing $\bm{R}$ as a sequence $\bm{B}_m\bm{B}_{m-1}\cdots\bm{B}_1$ can be seen as facilitating multi-stage information exchanges, where each $\bm{B}_i$ specifies which connections exist at the $i$-th stage. The process starts with $\bm{B}_1$ and culminates with $\bm{B}_n$. As an illustration in Figure~\ref{fig:info_trans}, we examine the decomposition $\bm{R}=\bm{B}_5\bm{B}_4\bm{B}_3\bm{B}_2\bm{B}_1$, showing both the individual sparsity patterns and the corresponding transmission graph. This graph arises by interpreting the product as a chain of connectivity: each $\bm{B}_i$ links nodes in the $i$-th layer to those in the $(i+1)$-th, with the entry $(j_1, j_2)$ of $\bm{B}_i$ indicating whether a directed connection exists from node $j_2$ at level $i$ to node $j_1$ at level $i+1$ (a zero indicating no link). To ensure that $\bm{B}_5\bm{B}_4\bm{B}_3\bm{B}_2\bm{B}_1$ forms a dense matrix, there must be paths allowing every node at level one to reach all nodes at level six. If we restrict attention to the product $\bm{R}=\bm{B}_4\bm{B}_3\bm{B}_2\bm{B}_1$, corresponding to information flow from the first layer to the fifth, we observe that node 1 at the source cannot reach node 3 at the target; thus, the resulting matrix has a zero at position $(3,1)$. This relationship between reachability in the induced graph and zeros in the product persists for $\bm{R}=\bm{B}_3\bm{B}_2\bm{B}_1$ and $\bm{R}=\bm{B}_2\bm{B}_1$ as well. More broadly, for $\bm{R}\in\mathbb{R}^{d\times d}$ to be fully dense, the factors $\bm{B}_m,\cdots,\bm{B}_1$ must define directed edges in a $d\times (m+1)$ grid, where each edge joins only adjacent layers (that is, columns), guaranteeing that each node in the initial layer can transmit information to all nodes in the terminal layer.

The information transmission perspective offers valuable insights into the sparsity characteristics observed in the factorization matrices utilized in OFT. As depicted in Figure~\ref{fig:blk_diag}, the original OFT displays a block-diagonal structure for $\bm{R}$. While this configuration decreases the number of parameters subject to training, it inherently lacks the ability to form a densely populated $\bm{R}$ matrix. Our primary objective is to synthesize a dense orthogonal matrix by combining $m$ sparse orthogonal matrices, all while minimizing the count of effective trainable parameters. Within the context of the information transmission approach, two principal criteria guide this objective: \emph{(i)} \textbf{dense connectivity}, where each node at the initial layer maintains at least one connection pathway to every node at the terminal layer, and \emph{(ii)} \textbf{minimum free edges}, demanding the total edge count be minimized under the constraint of orthogonality. 

Orthogonality introduces nuanced restrictions regarding edge placement between consecutive layers. Specifically, to guarantee that each $\bm{B}_i$ matrix is full-rank—a prerequisite for orthogonality—it is essential to establish $d$ edges that create a bijective mapping between the nodes of the $i$-th layer and those of the $(i-1)$-th layer. This requirement ensures that adjacent layers have no fewer than $d$ interconnections (\emph{e.g.}, as shown with $d = 4$ in Figure~\ref{fig:info_trans}). These $d$ connections are indispensable for preserving orthogonality and should not contribute to the effective trainable edge count, as the associated matrix elements are fixed and non-trainable (\emph{e.g.}, a $d\times d$ orthogonal matrix with $d$ non-zero entries, each constrained to $\pm 1$). Furthermore, because orthogonal matrices are parameter-efficient compared to dense ones, the orthogonality condition induces dependencies between the edges. For instance, in a $2\times 2$ orthogonal matrix, if the $(1,1)$ entry is zero, orthogonality compels the $(2,2)$ entry to also be zero (\emph{i.e.}, the absence of one edge necessitates the absence of another). Consequently, for any set of permissible edge arrangements, orthogonality may result in the addition or removal of specific edges. By examining the non-zero structure in the resultant orthogonal matrix, the information transmission framework proves particularly advantageous in the OFT setting, as our primary interest lies in the configuration of non-zero, trainable elements in $\bm{R}$, rather than their exact numerical values.

A straightforward dense connection linking two layers requires $\mathcal{O}(d^2)$ edges, which corresponds to utilizing a full dense orthogonal matrix. This results in $d^2-d$ total edges; for instance, when $d=4$, this amounts to 12 edges. As shown in Figure~\ref{fig:info_trans}, an example of a possible matrix factorization uses just $10$ edges, fewer than what a single dense orthogonal matrix would require. Such a viewpoint allows us to analyze OFT’s parameter efficiency using graph-theoretic intuition and systematically construct feasible factorizations. Our methodology is inspired by a notable network structure employed in the Cooley-Tukey algorithm, commonly referred to as butterfly graphs~\cite{cooley1965algorithm}. This architecture provides a dense connectivity pattern between $d$ input and $d$ output nodes using only $\mathcal{O}(d\log d)$ edges. For instance, while the configuration in Figure~\ref{fig:info_trans} achieves full connectivity with $10$ edges, the butterfly graph achieves the same with only $8$ edges. We proceed to explain how incorporating the butterfly structure leads to enhanced parameter efficiency.

\section{Orthogonal Parameterization by Butterfly Factorization}

The butterfly topology has its origins in the Cooley-Tukey algorithm, where it plays a key role in the fast Fourier transform. The essential idea is that in a Fourier transform, modifying a single frequency component can affect all points in the spatial domain, mirroring the problem of information transmission—each node at the initial layer can potentially transfer information to any node at the output layer. The butterfly configuration has also found widespread adoption as a network topology in computer architecture~\cite{solihin2015fundamentals,li2011linear} because of its efficacy for information distribution. Let $k\geq 2$ be a power of $2$. We first define the \emph{butterfly factor} $\bm{B}^F(k)$ as follows:

\begin{equation}\label{eq:but_factor}
\begin{aligned}
    \bm{B}^F(k)=\begin{bmatrix}
    \text{diag}(\bm{d}_1) & \text{diag}(\bm{d}_2) \\
    \text{diag}(\bm{d}_3) & \text{diag}(\bm{d}_4)
    \end{bmatrix}\in\mathbb{R}^{k\times k},
\end{aligned}
\end{equation}

with each $\bm{d}_i\in\mathbb{R}^{\frac{k}{2}}$ for $i=1,\dots,4$ being suitable vectors. Next, assuming $d=2^N$, we define a $d$-dimensional \emph{butterfly matrix} $\bm{B}(d)\in\mathbb{R}^{d\times d}$ recursively in the following way:

\begin{equation}
\begin{aligned}
    \!\!\bm{B}(d)=\tilde{\bm{B}}(d,d)\!\cdot\!\begin{bmatrix}
    \bm{B}_1(\frac{d}{2})\!\!\!\!\! & \bm{0} \\
    \bm{0}\!\!\!\!\! &\bm{B}_2(\frac{d}{2})
    \end{bmatrix}= \tilde{\bm{B}}(d,d)\tilde{\bm{B}}(d,\frac{d}{2})\cdots\tilde{\bm{B}}(d,2),
\end{aligned}
\end{equation}

where $\bm{B}_1(\frac{d}{2})$ and $\bm{B}_2(\frac{d}{2})$ represent two butterfly matrices, each of dimension $\frac{d}{2}$. The \emph{butterfly component} is formulated as $\thickmuskip=2mu \medmuskip=2mu \tilde{\bm{B}}(d,k)=\text{diag}(\bm{B}^F_1(k),\cdots,\bm{B}^F_{d/k}(k))$, yielding a $d\times d$ block-diagonal matrix with blocks of size $k$, where each $\bm{B}^F_i(k)$ is a butterfly factor defined as in Equation~\ref{eq:but_factor}. With this setup, butterfly matrices are employed to represent orthogonal matrices. The key requirement is to ensure that every multiplicative factor $\tilde{\bm{B}}(d,k)$ comprising the butterfly matrix $\bm{B}(d)$ is orthogonal. Examining the case when $k=2$, the block-diagonal matrix $\tilde{\bm{B}}(d,2)$ can be directly parameterized to be orthogonal using the Cayley transform (or, equivalently, a $2$-dimensional rotation), as detailed by \cite{liu2021orthogonal,qiu2023controlling}. The arrangement of non-zero entries in each butterfly component aligns, up to permutation, with that of $\tilde{\bm{B}}(d,2)$, making the orthogonal parameterization straightforward for all butterfly components. This structure provides a computationally efficient scheme for parameterizing orthogonal matrices using collections of $2\times 2$ orthogonal blocks. Building on \cite{chen2022pixelated}, this approach can be extended to the block setting by introducing the block butterfly component $\tilde{\bm{B}}^b(d,k)$, where each element $\bm{d}_i$ is now a $b\times b$ matrix. For $\tilde{\bm{B}}^b(d,2)$ to maintain orthogonality, each $2b\times 2b$ block must be parameterized as an orthogonal matrix. Similarly, for $k>2$, the support pattern of non-zero blocks in $\tilde{\bm{B}}^b(d,k)$ can be obtained by block-wise permutation of that in $\tilde{\bm{B}}^b(d,2)$, allowing the same orthogonal parameterization. By combining these elements, the forward computation in BOFT becomes

\begin{equation}
    \begin{aligned}\nonumber
    \bm{z} = \big(\bm{R}(m,b) \cdot \bm{W}^0\big)^\top \bm{x},~~~~\text{s.t.}~\left\{ \bm{R}(m,b) = \prod_{i=1}^{m} \tilde{\bm{B}}^{b}_{(d,i)}~~\&~~\underbrace{\big(\tilde{\bm{B}}^b_{(d,j)}\big)^\top\tilde{\bm{B}}^b_{(d,j)} = \tilde{\bm{B}}^b_{(d,j)} \big(\tilde{\bm{B}}^b_{(d,j)}\big)^\top\! = \bm{I}_d}_{\forall j \in [1, m]} \right\},
    \end{aligned}
\end{equation}

Here, for the sake of brevity, we use $\tilde{\bm{B}}^b(d,2^{m-i+1})$ to denote $\tilde{\bm{B}}^b_{(d,i)}$, and $\bm{I}_d$ refers to the $d$-dimensional identity matrix. The matrix $\bm{R}(m,b) \in \mathbb{R}^{d \times d}$ is a product of $m$ orthogonal butterfly factors, resulting in an orthogonal matrix. For succinct notation, we refer to BOFT instantiated with $\bm{R}(m, \frac{b}{2})$ as BOFT($m$, $b$), with $b \geq 2$. In the case when $m=1$, BOFT($1$,$b$) simplifies to the block-diagonal OFT~\cite{qiu2023controlling}, where the block size is $b$. For $b = d$, BOFT($1$, $d$) coincides with the standard OFT~\cite{qiu2023controlling} utilizing a fully unconstrained orthogonal matrix. When we set $m = \log \frac{2d}{b}$, BOFT($\log \frac{2d}{b}$, $b$) applies the block butterfly matrix $\bm{B}^{\frac{b}{2}}(d)$ as $\bm{R}$, giving rise to a full, dense orthogonal matrix. Broadly, BOFT($m$,$b$) parameterizes $\frac{1}{2}(b-1)dm$ effective finetunable parameters for reparameterizing a $d \times n$ linear transformation. Employing the standard butterfly construction, i.e., $m = \log d, b = 2$, leads to $\mathcal{O}(d \log d)$ parameters in BOFT. By comparison, the classic OFT, which uses an unrestricted orthogonal matrix, incurs $\mathcal{O}(d^2)$ parameters, and the block-diagonal OFT with $r$ blocks has $\mathcal{O}(bd)$ parameters. Thus, the dense orthogonal parameterization in the original OFT can only be attained by taking $b=d$, while in BOFT, dense orthogonal matrices can be realized for arbitrary values of $b$.

\textbf{Identity Initialization in BOFT}. Common finetuning approaches typically initialize from the parameters of the pretrained network to minimize deviation from the pretrained performance. Analogous to LoRA, which sets low-rank weights initially to zero, BOFT utilizes identity matrices to initialize all butterfly components (meaning all skew-symmetric matrices in the Cayley transform are zero-initialized).

\textbf{Multiplicative Dropout in BOFT}. LoRA~\cite{hulora2022} incorporates a Dropout module for its low-rank parameter updates to help prevent overfitting, where classic Dropout~\cite{srivastava2014dropout} can be used as-is. However, standard Dropout is incompatible with BOFT due to the multiplicative formulation of the orthogonal weight update. To resolve this, we devise a multiplicative Dropout for BOFT. Given that $\bm{R}(m,b)$ consists of $m$ butterfly-structured orthogonal matrices, each of which can be reorganized into $2b \times 2b$ block-diagonal forms, multiplicative Dropout is realized by randomly selecting a fraction $p_1$ of butterfly components and a fraction $p_2$ of diagonal blocks within each component to be replaced with identity blocks.

\section{Intriguing Insights and Discussions}

\textbf{Expressivity of BOFT}. The combination of the butterfly structure and permutations allows for the exact reconstruction of numerous well-known fast linear transforms~\cite{dao2019learning,dao2019kaleidoscope} (\emph{e.g.}, fast Fourier transform, Hadamard transform). Nevertheless, the capability of our orthogonal butterfly matrix to closely approximate an arbitrary orthogonal matrix remains an open question. To explore this, we carry out a simulation in which a randomly generated dense orthogonal matrix~\cite{anderson1987generation} of size $512\times 512$ is approximated. The outcomes, averaged over 10 random seeds, are displayed in Figure~\ref{fig:para_eff}. On these plots, the y-axis represents the approximation error, and the x-axis represents the count of effective trainable parameters. For each color, the curves correspond to BOFTs with fixed block sizes; the point at the far left corresponds to BOFT($1$,$b$) (\emph{i.e.}, the conventional block-diagonal OFT with block size $b$). Across the board, BOFT demonstrates superior parameter efficiency in comparison with OFT. For instance, BOFT($9$,$2$) exhibits higher expressivity than BOFT($1$,$16$), despite having significantly fewer parameters. Configurations with smaller $b$ and larger $m$ tend to be more parameter-efficient; for example, BOFT($6$,$4$) utilizes considerably fewer parameters while achieving approximation errors similar to those of BOFT($2$,$16$). Overall, the butterfly matrix parameterizes a more restricted, yet highly structured, subset within the orthogonal group relative to the block-diagonal form, rendering BOFT provably more expressive than OFT at a fixed block size.

\begin{theorem}[Expressivity of BOFT]\label{thm:exp_boft}
    BOFT possesses greater expressivity than OFT given an identical block size. To enable the butterfly matrix to approximate any orthogonal matrix of size $d$, one can form products such as $\bm{B}_{d-1,1}(d)\bm{B}^\top_{d-1,2}(d)\cdots\bm{B}_{1,1}(d)\bm{B}^\top_{1,2}(d)$, where all $\bm{B}_{i,j}(d)$ denote butterfly matrices for all $i$ and $j$.
\end{theorem}

Theorem~\ref{thm:exp_boft} points to a natural extension for BOFT—constructing the final orthogonal matrix as $\thickmuskip=2mu \medmuskip=2mu \bm{R}^G(m_1,b_1,m_2,b_2,l)=\bm{R}_{l,1}(m_1,b_1)\bm{R}_{l,2}^T(m_2,b_2)\cdots\bm{R}_{1,1}(m_1,b_1)\bm{R}_{1,2}^T(m_2,b_2)$, where $\bm{R}_{i,j}^T(m,b)$ represents an orthogonal matrix in BOFT. When taking $m_1=m_2=\log d$, $b_1=b_2=2$, and $l=d-1$, this construction $\bm{R}^G(m_1,b_1,m_2,b_2,l)$ is able to span the full orthogonal group. This manner of composing matrices aligns with what is termed the kaleidoscope hierarchy~\cite{dao2019kaleidoscope}. It should be emphasized, however, that an increase in expressivity does not automatically result in better outcomes when fine-tuning: full fine-tuning, though universally expressive, frequently delivers suboptimal performance. Ultimately, the interplay between expressivity and structural regularization determines the ability of a model to generalize after fine-tuning. By expanding the parameter search space through structural priors, BOFT facilitates a more favorable balance between expressivity and regularity.

\textbf{Spectral properties}. Orthogonal finetuning consistently achieves superior spectral characteristics compared to LoRA, as it fully maintains the spectral norm of the original weight matrix $\bm{W}^0$. To demonstrate this, consider the singular value decomposition $\bm{W}^0=\bm{U}\bm{\Sigma}\bm{V}^\top$, where $\bm{U}$ and $\bm{V}$ are orthogonal matrices, and $\bm{\Sigma}$ contains the singular values on its diagonal. The procedures OFT and BOFT both involve left-multiplying by an orthogonal matrix $\bm{R}$, resulting in a finetuned weight of $\bm{R}\bm{U}\bm{\Sigma}\bm{V}^\top$. This operation leaves the top singular value, and thus the spectral norm of $\bm{W}^0$, unchanged. Maintaining this property has been empirically shown to enhance both training stability and the model's generalization abilities~\cite{miyato2018spectral,yoshida2017spectral}. Further mathematical aspects are detailed in Appendix~\ref{app:sn_property}.

\textbf{Orthogonal finetuning as learning bilinear similarity}. BOFT can also be interpreted as learning a bilinear similarity function of the form $\bm{w}^0_i\bm{R}\bm{x}$, where $\bm{w}^0_i$ denotes the $i$-th column (or neuron) of $\bm{W}^0$. In this context, BOFT is learning the bilinear similarity matrix $\bm{R}$ subject to orthogonality constraints, closely relating it to distance metric learning~\cite{xing2002distance} and classical bilinear forms~\cite{roman2005advanced}.

\textbf{Inductive bias and generalization in BOFT}. The matrix $\bm{R}(m,b)$ in BOFT typically lies within a structured subset of the orthogonal group, restricting the hypothesis space and thereby inducing a specific inductive bias. We contend that this form of bias—stemming from the butterfly factorization—positively impacts generalization, as it encodes structural patterns present in numerous standard linear transformations~\cite{dao2019kaleidoscope} such as the discrete Fourier transform, discrete sine/cosine transform, and the Hadamard transform. In addition, the utilization of sparse matrix factorizations in BOFT can also introduce implicit forms of inductive bias~\cite{gunasekar2017implicit,li2018algorithmic,liu2019neural}.

\textbf{Relation to butterfly-based sparse training}. A number of prior works~\cite{dao2019learning,dao2019kaleidoscope,chen2022pixelated,dao2022monarch} investigate the application of butterfly parameterization within sparse training regimes. These studies generally emphasize direct reparameterization of weight matrices using the butterfly structure, initializing and training neural networks from the beginning. Notably, \cite{dao2022monarch} explores finetuning pretrained weights by initially projecting them onto modified butterfly matrices, followed by optimization of these projected segments for specific downstream applications. In contrast, BOFT introduces a novel finetuning method that modifies weights using layer-shared weight matrices, diverging fundamentally from existing strategies.

\input{rebuttal-results}

\section{Concluding Remarks and Limitations}

In this work, we present Orthogonal Butterfly, a broadly applicable and parameter-efficient finetuning approach for foundation models grounded in the butterfly architecture. The primary conceptual advancement lies in achieving higher parameter-efficiency by representing a dense orthogonal matrix as a composition of several sparse orthogonal matrices. To systematically identify valid matrix decompositions, we introduce a graphical framework based on information transmission. This perspective reveals that the butterfly structure can satisfy our requirements for sparse orthogonal matrix factorization particularly well. Through extensive experiments, we validate the practical utility of BOFT when finetuning models in large-scale language, vision, and text-to-image generative contexts. Empirical findings also highlight BOFT's advantages as a universal finetuning solution.

Nevertheless, BOFT does have limitations. Specifically, since the resultant orthogonal matrix in BOFT is formed by multiplying several orthogonal matrices, this construction incurs a slightly greater training computational cost relative to OFT. Addressing the issue of BOFT’s training overhead remains unresolved. On the positive side, after finetuning, the orthogonal matrices derived via BOFT can be merged into the original model weights, leading to no additional cost at inference time. Furthermore, it remains unclear whether the butterfly network provides the optimal structure for information transmission. Our proposed framework offers opportunities to leverage insights from other disciplines such as computer networking, where optimality of network topology for information dissemination has been rigorously analyzed. Future work may uncover network architectures superior to the butterfly design for assembling dense orthogonal matrices.